\chapter{Creio em Deus Pai todo-poderoso, criador do céu e da terra}
\chaptermark{Creio em Deus Pai todo-poderoso}

O primeiro artigo da Profissão de Fé, ``Creio em Deus Pai todo-poderoso, criador do céu e da terra'', forma o fundamento de toda a doutrina e da vida cristã. Ele sintetiza a revelação de Deus como único, apresentado simultaneamente como Pai, como Ser dotado de poder infinito e como Criador de toda a realidade visível e invisível. Este artigo não apenas inicia o Credo, mas sustenta todos os demais mistérios da fé, revelando a relação entre Deus e a criação, entre o Criador e o homem, e entre a bondade original e a realidade do pecado. O texto a seguir expõe, de forma ordenada e ampliada, os principais ensinamentos do Catecismo da Igreja Católica sobre cada elemento deste artigo, incluindo o relato do livro do Gênesis desde a criação até a aliança com Noé, com particular atenção à degeneração progressiva da humanidade após o pecado, às prefigurações de Jesus Cristo e da Igreja presentes nestas páginas e à promessa da redenção que aponta para a obra salvífica de Cristo.

\section{Creio em Deus (199--231)}

A confissão ``Creio em Deus'' constitui a afirmação mais fundamental de todo o Credo Apostólico. Todo o conteúdo da fé cristã depende desta primeira e radical verdade: existe um só Deus, e é sempre em relação a Ele que se compreendem o homem, o mundo e todos os outros artigos da fé. Esta profissão de fé não é apenas o ponto de partida, mas o alicerce que sustenta toda a vida do cristão, orientando-o para Aquele que é o princípio e o fim de todas as coisas.

A fé no Deus único tem as suas raízes na revelação feita ao povo de Israel e manifesta-se como uma profissão inseparável da própria existência da fé cristã. Deus revelou-se progressivamente como o Senhor único e, através dos profetas, estende a todos os povos da terra o convite à salvação, proclamando que só n’Ele se encontram a justiça e a força. Esta certeza oferece à existência cristã um alicerce sólido e permanente, permitindo compreender que toda a vida do fiel deve orientar-se para este Deus único.

Esta confissão distingue claramente a fé cristã de outras concepções religiosas, estabelecendo desde o início que a relação com Deus não é uma entre muitas possibilidades, mas a única realidade fundamental sobre a qual se constrói toda a experiência de salvação. A verdade fundamental de que Deus é o Primeiro e o Último prepara o terreno para a compreensão mais profunda do mistério trinitário e da obra criadora de Deus. Deste modo, o crente é convidado a colocar toda a sua confiança n’Aquele que transcende o mundo e a história, mas que ao mesmo tempo se aproxima do homem com amor.

Ao professar ``Creio em Deus'', o cristão reconhece que Deus é o fundamento de toda a realidade criada. Esta fé ilumina a inteligência e fortalece a vontade para viver de acordo com a dignidade de quem foi criado à imagem e semelhança de Deus. É uma profissão que transforma toda a existência, chamando o fiel a viver em constante referência a este Deus único, bom e misericordioso.

\section{O Pai (232--267)}

Os cristãos recebem o Batismo ``em nome do Pai e do Filho e do Espírito Santo''. Antes de receber este sacramento, o fiel responde afirmativamente à profissão de fé em cada uma das três Pessoas divinas. Existe apenas um Deus: o Pai todo-poderoso, o seu Filho único e o Espírito Santo. Esta verdade central revela que a paternidade de Deus não é uma imagem humana aplicada a Deus, mas a própria realidade divina da qual deriva toda a paternidade na terra.

O mistério da Santíssima Trindade é o mistério central de toda a fé e de toda a vida cristã, constituindo a fonte de todos os outros mistérios e a luz que os ilumina. Esta verdade não se apresenta como uma questão meramente teórica ou especulativa, mas como a realidade mais íntima na qual o cristão é introduzido pelos sacramentos da iniciação cristã. A confissão de Deus como Pai abre o caminho para compreender a paternidade divina como origem de toda a vida e de toda a graça.

A dimensão trinitária da fé permite uma compreensão integrada de todos os artigos subsequentes do Credo, revelando que a criação, a salvação e a santificação têm a sua origem e o seu termo na comunhão eterna do Pai, do Filho e do Espírito Santo. Deus revela-se como Pai amoroso que deseja comunicar a sua própria vida aos homens. Esta paternidade divina convida cada pessoa a entrar numa relação filial de confiança e amor.

Ao reconhecer Deus como Pai, o cristão descobre a própria identidade como filho amado. Esta verdade transforma a oração, a vida moral e a compreensão da Igreja como família de Deus. É uma revelação que enche de esperança o coração humano, pois garante que nunca estamos sós, mas sempre acompanhados pelo amor do Pai celestial.

A profissão de fé em Deus Pai convida ainda à conversão do coração, para que o fiel viva como verdadeiro filho, imitando a obediência e o amor do Filho Jesus. Desta forma, o primeiro artigo do Credo torna-se caminho de vida e fonte de alegria para o cristão.

\section{O Todo-Poderoso (268--278)}

De todos os atributos divinos, apenas a onipotência de Deus é explicitamente nomeada no Credo, o que revela a importância central desta verdade para a vida de fé. Acredita-se que o poder de Deus é universal, porque Ele criou todas as coisas, as governa e pode fazer tudo o que deseja. Este poder não é arbitrário, mas sempre exercido com sabedoria e amor.

Este poder é ao mesmo tempo amoroso, pois Deus age como Pai providente, e misterioso, pois somente a fé consegue discerni-lo, especialmente quando se manifesta perfeito na fraqueza. As Sagradas Escrituras proclamam repetidamente que nada é impossível para Deus, que dispõe todas as obras segundo a sua vontade e que o seu domínio se estende sobre o céu e sobre a terra. A onipotência divina revela-se sobretudo no cuidado providente pelas necessidades humanas.

A onipotência de Deus manifesta-se de modo privilegiado na misericórdia ao perdoar os pecados e na adoção filial que oferece aos homens. Longe de ser um poder distante ou temível, a onipotência de Deus Pai é força de amor que acompanha a história humana, sustenta a criação e realiza a salvação. Esta verdade oferece ao fiel a confiança necessária para enfrentar as dificuldades da existência.

A omnipotência divina ilumina igualmente a compreensão do mistério do mal e do sofrimento. Embora Deus permita o livre arbítrio humano, o seu poder amoroso é capaz de tirar bem até do mal. O cristão encontra nesta verdade motivo de esperança e de abandono confiante nas mãos do Pai.

Finalmente, confessar Deus como Todo-Poderoso significa reconhecer que a história está nas suas mãos e que o seu plano de amor se cumprirá plenamente. Esta fé dá coragem para viver com generosidade e perseverança no caminho da santidade.

\section{O Criador (279--324)}

A profissão de fé inicia com as palavras solenes ``No princípio Deus criou o céu e a terra''. A criação representa o fundamento de todos os desígnios salvíficos de Deus e marca o verdadeiro início da história da salvação, que encontra o seu pleno cumprimento em Jesus Cristo. O mistério de Cristo ilumina o mistério da criação e revela o fim último para o qual Deus criou todas as coisas.

A criação não deve ser vista como um tema isolado ou meramente cosmológico, mas como a base indispensável sobre a qual se compreende toda a obra salvífica de Deus, que culmina nos sacramentos da Igreja. Por este motivo, tanto a liturgia da Vigília Pascal como a catequese antiga destinada aos que se preparavam para o Batismo iniciavam sempre pelo relato da criação. A doutrina da criação revela um Deus que age por amor e com total liberdade.

Deus criou o mundo pela sua Palavra, manifestando a sua sabedoria e bondade. Cada criatura reflete, a seu modo, a beleza e a perfeição do Criador. O ato criador de Deus continua a sustentar a existência de todas as coisas a cada momento.

Esta verdade convida o cristão a contemplar a natureza com respeito e gratidão, reconhecendo nela o dom do Pai. A criação é um livro aberto que revela a glória de Deus àqueles que têm olhos de fé. Ao mesmo tempo, ela nos recorda a nossa responsabilidade de cuidar da obra das mãos de Deus.

\section{O Céu e a Terra (325--354)}

O Credo professa que Deus é ``criador do céu e da terra'', expressão que abrange tudo o que existe, tanto o visível como o invisível. A fórmula ``céu e terra'' designa a totalidade da criação e indica ao mesmo tempo a distinção e a profunda unidade que existe entre o mundo material e o mundo espiritual. Esta verdade abarca tanto os anjos como o universo visível.

O Quarto Concílio do Latrão declara solenemente que Deus, desde o princípio do tempo, criou simultaneamente as duas ordens de criaturas — a espiritual e a corporal — e, em seguida, o homem, que participa de ambas as ordens. Esta doutrina ajuda a perceber a unidade e a harmonia do projeto divino sobre toda a realidade. O céu e a terra proclamam a glória do Criador.

A distinção entre o visível e o invisível não implica oposição, mas complementaridade no plano de Deus. O mundo material é bom e ordenado para o serviço do homem e para a glória de Deus. Os anjos, por sua vez, são mensageiros e servos do Altíssimo.

Esta visão da criação convida o cristão a viver em harmonia com todo o universo criado, louvando o Criador por meio das coisas visíveis e invisíveis. A profissão de fé no Criador do céu e da terra fortalece a esperança na vida eterna, onde o céu e a terra se unirão na nova criação.

\section{O homem (355--384)}

O ser humano ocupa um lugar único e singular em toda a criação, tendo sido criado ``à imagem de Deus'', homem e mulher. Ele une em si mesmo o mundo espiritual e o mundo material, foi estabelecido na amizade com Deus e é o único ser visível capaz de conhecer e amar o seu Criador. Esta dignidade inviolável fundamenta-se no amor incalculável de Deus.

Sendo imagem de Deus, o homem possui a dignidade de pessoa: é capaz de autoconhecimento, de autodeterminação e de se dar livremente em comunhão com os outros e com Deus. Esta visão da pessoa humana oferece a base antropológica indispensável para compreender o sentido profundo dos sacramentos da iniciação cristã. O homem foi criado para a felicidade eterna.

A dignidade do homem permanece mesmo depois do pecado, embora ferida. Deus nunca deixou de amar a criatura humana e preparou desde o início o caminho da salvação. Esta verdade deve levar-nos a respeitar a vida humana em todas as suas etapas.

O ser humano encontra a sua realização plena somente em Deus. A sua vocação é viver em comunhão com o Criador e com os irmãos. Esta chamada ressoa no coração de cada pessoa como sede de infinito.

\section{A queda (385--421)}

A doutrina da criação inclui necessariamente o relato da queda no pecado, da qual Jesus Cristo veio libertar a humanidade. Depois de apresentar o Criador e a bondade original da criação, a fé contempla o drama do pecado original e as suas graves consequências para a história humana. O pecado não destruiu definitivamente o projeto amoroso de Deus.

O pecado original feriu a natureza humana, introduzindo a inclinação para o mal e a desarmonia nas relações. No entanto, Deus prometeu imediatamente um Redentor, mostrando que a sua misericórdia é maior que o pecado. Esta promessa percorre toda a história da salvação.

A queda revela a seriedade da liberdade humana e as consequências do afastamento de Deus. Ao mesmo tempo, abre o espaço para a manifestação do amor misericordioso do Pai que envia o seu Filho para salvar o que estava perdido.

A Igreja ensina que todos os homens nascem afetados pelo pecado original e necessitam da graça do Batismo. Esta verdade ajuda a compreender a condição humana e a necessidade da redenção em Cristo.

\section{O Livro do Gênesis: Da Criação até a Aliança com Noé}

O livro do Gênesis narra que Deus criou o universo inteiro pela sua Palavra poderosa, separando a luz das trevas, fazendo surgir a vida em todas as suas formas e culminando esta obra na criação do ser humano à sua imagem e semelhança. Colocou o homem e a mulher no jardim do Éden e estabeleceu com eles uma relação íntima de amizade e confiança. A criação é declarada boa por Deus em cada etapa.

Após o pecado original, a humanidade inicia um processo visível e progressivo de degeneração moral e espiritual. A desobediência de Adão e Eva introduz imediatamente a vergonha, o medo e a divisão entre eles. Caim assassina o seu irmão Abel, demonstrando como o pecado gera violência e morte. A maldade alastra-se de tal forma que culmina no dilúvio como ato de purificação e justiça divina.

No meio desta degradação generalizada, Deus escolhe Noé, homem justo, e manda-o construir a arca que se torna instrumento de salvação para a sua família e para as espécies animais. Após o dilúvio, Noé oferece um sacrifício agradável a Deus, que estabelece então uma aliança com toda a humanidade, prometendo nunca mais destruir a terra com as águas e colocando o arco-íris no céu como sinal perpétuo da sua fidelidade.

\section{Prefigurações de Jesus Cristo e da Igreja}

No relato do Gênesis encontram-se várias prefigurações significativas de Jesus Cristo e da Igreja. Adão aparece claramente como figura de Cristo, o novo Adão que viria reparar a desobediência do primeiro. A criação da mulher a partir da costela de Adão prefigura a Igreja nascida do lado aberto de Cristo na cruz.

A arca de Noé constitui uma imagem particularmente clara da Igreja: fora dela não há salvação, e dentro dela os fiéis são preservados e salvos pelas águas que, no Novo Testamento, se tornam as águas regeneradoras do Batismo. Estas prefigurações revelam a unidade do plano de Deus ao longo da história.

Deus age de forma progressiva, preparando a humanidade para a plenitude da revelação em Cristo. Cada episódio do Gênesis aponta misteriosamente para o mistério pascal.

\section{A Prefiguração da Redenção da humanidade por Jesus Cristo}

Toda a narrativa do Gênesis demonstra de modo eloquente que o pecado, por mais grave que seja, não tem a última palavra na história humana. A salvação da humanidade através das águas do dilúvio aponta de forma clara para o Batismo, pelo qual Cristo nos liberta do pecado. A aliança estabelecida com Noé antecipa a Nova e Eterna Aliança que seria selada no sangue de Cristo.

A fidelidade de Deus manifestada no sinal do arco-íris encontra a sua plena realização na Ressurreição de Cristo, que representa a vitória definitiva sobre o pecado, a morte e toda a degeneração da humanidade, abrindo o caminho para a nova criação prometida por Deus desde o princípio. Esta promessa de redenção enche o coração do crente de esperança.

A história contada no Gênesis culmina na certeza de que Deus é fiel às suas promessas. O cristão é chamado a confiar nesta fidelidade e a viver desde já como filho da nova criação inaugurada por Jesus Cristo.